%----------------------------------------------------------------------
% Appendix: Merging UCDP and ACLED with meltt
%----------------------------------------------------------------------
\subsection{Merging the UCDP and ACLED event data}
\label{app:meltt}

We work with the UCDP Georeferenced Event Dataset (GED) and the ACLED
dataset, in both cases restricted to events located in Syria.\footnote{UCDP
GED, version downloaded on 6 February 2026; ACLED Syria export, downloaded on
4 February 2025.} The UCDP extract covers the period from 18 March 2011 to
31 December 2024 and contains $87{,}857$ events, of which $73{,}755$ are
state-based, $10{,}092$ non-state and $4{,}010$ one-sided events.\footnote{We
drop four isolated UCDP events that predate the Syrian civil war: three events
dated 2004 (one Israel--Hamas clash and two events of one-sided violence by the
Government of Syria against civilians) and one Israel--Hezbollah event recorded
in February 2008. All four are incidents that merely take place on Syrian
territory; UCDP coverage of the Syrian conflict itself is continuous only from
2011 onwards.} The ACLED extract covers the period from 1 January 2017 to
31 January 2025 and contains $28{,}241$
events, all of which are coded as \emph{Battles}.

To combine the two datasets we use the \texttt{R} package \texttt{meltt}
(Merging Event Data by Location, Time and Type; \citealp{Donnay2019}), which
integrates event data from different sources and removes duplicate reports of
the same incident. The merge is based on space, time, and two user-defined
taxonomies that render the source-specific event categories comparable. The
spatial-precision taxonomy maps UCDP's \texttt{where\_prec} (seven categories)
and ACLED's \texttt{geo\_precision} (three categories) onto common levels
(exact location vs.\ administrative unit vs.\ country; sub-regional vs.\
sub-national vs.\ national; precise vs.\ imprecise). The violence taxonomy maps
UCDP's \texttt{type\_of\_violence} and ACLED's \texttt{interaction} variable
onto the common categories \emph{state-based}, \emph{non-state} and
\emph{other}. Two events are treated as potential duplicates if they lie within
a temporal window of one day (\texttt{twindow}~$=1$) and a spatial window of
three kilometres (\texttt{spatwindow}~$=3$) of each other and are compatible
along the levels of both taxonomies.

The procedure identifies $5{,}198$ duplicate event pairs across the two
sources. After collapsing these, the merged dataset contains $110{,}899$ unique
events, of which $82{,}658$ originate from UCDP and $28{,}241$ from ACLED.
Table~\ref{tab:meltt} summarises the two input datasets and the merged output.

\begin{table}[htbp]
\centering
\caption{Input datasets and \texttt{meltt} output}
\label{tab:meltt}
\begin{tabular}{lccc}
\hline\hline
 & Period & Events & Contribution to merged data \\
\hline
UCDP GED  & 18 Mar 2011 -- 31 Dec 2024 & $87{,}857$ & $82{,}658$ \\
ACLED     & 1 Jan 2017 -- 31 Jan 2025  & $28{,}241$ & $28{,}241$ \\
\hline
Duplicates identified & \multicolumn{3}{c}{$5{,}198$ event pairs} \\
Merged dataset        & 18 Mar 2011 -- 31 Jan 2025 & \multicolumn{2}{c}{$110{,}899$} \\
\hline\hline
\end{tabular}
\end{table}

